% Теоретическая глава: почему исполнительную часть VLA-политики можно обучить
% только консервативными методами RL. Сборка: latexmk -xelatex main.tex
% Численная проверка неравенств: python3 verify_theorems.py
\documentclass[11pt,a4paper]{article}

\usepackage{fontspec}
\setmainfont{cmunrm.otf}[BoldFont=cmunbx.otf,ItalicFont=cmunti.otf,BoldItalicFont=cmunbi.otf]
\setsansfont{cmunss.otf}
\setmonofont{cmuntt.otf}[ItalicFont=cmunit.otf]
\usepackage[english,russian]{babel}
\usepackage[a4paper,margin=2.3cm]{geometry}
\usepackage{amsmath,amssymb,amsthm,mathtools,bm}
\usepackage{enumitem,booktabs,array,tabularx,longtable}
\usepackage[dvipsnames]{xcolor}
\usepackage[colorlinks=true,linkcolor=MidnightBlue,citecolor=MidnightBlue,urlcolor=MidnightBlue]{hyperref}

\setlist{itemsep=2pt,topsep=4pt}
\setlength{\parskip}{3pt}
\tolerance=2000
\emergencystretch=3em
\def\UrlBreaks{\do\/\do\_\do\-}

\theoremstyle{plain}
\newtheorem{theorem}{Теорема}[section]
\newtheorem{lemma}[theorem]{Лемма}
\newtheorem{proposition}[theorem]{Предложение}
\newtheorem{corollary}[theorem]{Следствие}
\theoremstyle{definition}
\newtheorem{definition}[theorem]{Определение}
\newtheorem{assumption}[theorem]{Допущение}
\theoremstyle{remark}
\newtheorem{remark}[theorem]{Замечание}

\newcommand{\E}{\mathbb{E}}
\newcommand{\PP}{\mathbb{P}}
\newcommand{\R}{\mathbb{R}}
\DeclareMathOperator{\TV}{TV}
\DeclareMathOperator{\KL}{KL}
\DeclareMathOperator{\clip}{clip}
\DeclareMathOperator*{\argmax}{arg\,max}
\newcommand{\cT}{\mathcal{T}}
\newcommand{\cZ}{\mathcal{Z}}
\newcommand{\cL}{\mathcal{L}}
\newcommand{\cD}{\mathcal{D}}
\newcommand{\cA}{\mathcal{A}}
\newcommand{\Pmix}{P_{\mathrm{mix}}}
\newcommand{\ab}{\bar\alpha}
\newcommand{\ip}[2]{\langle #1,#2\rangle}
\newcommand{\pit}{\tilde\pi}
\newcommand{\one}{\mathbf 1}

\title{Почему исполнительную часть VLA-политики\\можно обучить только обучением с подкреплением:\\математическое обоснование поэтапного бутстрэппинга}
\author{Д.\,Нуртдинов\\[2pt]\small Черновик теоретической главы диссертации}
\date{4 ноября 2026 г.}

\begin{document}
\maketitle

\begin{abstract}
Рассматривается адаптация предобученной модели «зрение--язык--действие» (VLA, здесь SmolVLA) к новой кинематической структуре (тросовый параллельный робот CDPR, манипулятор WidowX-200), при которой визуально-языковая часть заморожена, а обучается только исполнительная часть (остаточная голова действий) консервативными методами PPO/GRPO с усечением отношения правдоподобий. Ставится вопрос: почему такое обучение вообще возможно и приводит к успешному выполнению инструкций манипуляции. Ответ дан в виде системы утверждений с доказательствами. (1)~При бинарной награде GRPO в среднем поднимается по монотонному преобразованию $2\arcsin\sqrt{p}$ вероятности успеха, поэтому оптимизируется сама метрика задачи. (2)~Для вероятности успеха на конечном горизонте получена оценка улучшения в доверительной области и оценка сохранения навыков; из неё следуют условия, при которых навык не забывается, и объяснение накопленного дрейфа. (3)~Последовательность консервативных шагов трактуется как неточный зеркальный спуск в пространстве политик, который сходится к лучшей политике, представимой на замороженных признаках, со скоростью $O(\sqrt{H D/K})$. Здесь $D$ есть KL-расстояние от начальной политики до успешной. (4)~При разреженной награде сигнал отсутствует, пока вероятность успеха равна нулю, а для составной задачи она экспоненциально мала по числу этапов. Простые навыки, плотная награда, лестница стартовых распределений и композиция под одной инструкцией образуют конструктивный способ сделать её отделённой от нуля за полиномиальную цену. Сводная теорема объединяет эти утверждения. В ней точно определено, какую роль в оценках играли бы демонстрации и чем пайплайн их заменяет. Все неравенства проверены численно на точно решаемых конечных марковских процессах (приложение~\ref{app:numerics}).
\end{abstract}

\tableofcontents

%======================================================================
\section{Введение и постановка вопроса}\label{sec:intro}

\paragraph{Эмпирическая основа.} В эксперименте на тросовом роботе (CDPR) остаточный контроллер над замороженным SmolVLA достиг 69{,}14\,\% строгого успеха полной задачи \emph{put $\langle$object$\rangle$ into $\langle$destination$\rangle$} (1416 из 2048 попыток вне валидационной панели). Каждая попытка начиналась с пустого захвата и выполнялась без учителя. Поэтапный пайплайн на WidowX-200 воспроизводит ту же схему на другой кинематике: сначала простые навыки (подход, захват, размещение), затем их композиция под одной инструкцией.

\paragraph{Вопрос.} \emph{Почему исполнительную часть VLA-политики (action expert / actor / action head) можно обучить, используя только алгоритмы обучения с подкреплением с консервативными, слабо изменяющими политику обновлениями (PPO, GRPO с усечением), и почему это обучение приводит к успеху на инструкциях манипуляции объектами?}

Вопрос нетривиален по четырём причинам.
\begin{enumerate}[label=(\roman*)]
\item \textbf{Разреженность.} Награда бинарна: успех или неуспех. Пока успехов нет, градиент политики тождественно равен нулю (лемма~\ref{lem:stationary}).
\item \textbf{Длинный горизонт и контакт.} Успех требует цепочки точных событий: подход, захват, подъём, перенос, отпускание. Вероятность случайно пройти цепочку экспоненциально мала по её длине (предложение~\ref{prop:barrier}).
\item \textbf{Одна политика, много инструкций.} Параметры общие, поэтому улучшение одного навыка может разрушать другой. Это наблюдалось многократно (раздел~\ref{sec:experiments}).
\item \textbf{Новая кинематика.} Предобученные знания о действиях почти не переносятся: на CDPR замороженный SmolVLA-приор действует как постоянное смещение, а всё поведение задаёт остаточная голова. Поэтому из предобучения в задачу переносится прежде всего \emph{представление}, а не поведение.
\end{enumerate}

\paragraph{Что именно означает «только RL».} Чтобы утверждение было математическим, а не риторическим, зафиксируем информационную модель.

\begin{definition}[Информационная модель обучения]\label{def:info}
Обучающийся имеет доступ
(а) к замороженному предобученному отображению признаков $\varphi$ (визуально-языковой бэкбон, проекция, приор действий);
(б) к симулятору со сбросом в стартовые распределения из спецификации навыков (определение~\ref{def:skill});
(в) к предикатам успеха этих навыков.
Обучающийся \emph{не получает} целевых действий из внешнего источника: от человека, скриптового оракула или контроллера-помощника. Повторное использование данных, порождённых самой политикой (самоимитация на собственных успехах, переразметка инструкций), допускается: действия в этих данных выбраны обучающимся, а отобраны предикатом успеха.
\end{definition}

В этой модели «описать простейшие навыки» означает задать для каждого навыка тройку (инструкция, стартовое распределение, предикат успеха). Ниже показано, что через эти три объекта и замороженное представление выражаются все величины в оценках сходимости.

\paragraph{Структура ответа.} Ответ разбивается на четыре условия; каждому соответствует группа утверждений.
\begin{center}
\begin{tabularx}{\linewidth}{>{\raggedright}p{0.36\linewidth}X}
\toprule
Подвопрос & Результат \\
\midrule
(У1) Что оптимизирует GRPO? & Теорема~\ref{thm:arcsin}: монотонное преобразование вероятности успеха; следствие~\ref{cor:arcsin}: Парето-оптимальность. \\
(У2) Почему консервативный шаг не вредит? & Теорема~\ref{thm:improve} (улучшение), теорема~\ref{thm:retention} (сохранение навыков), следствия~\ref{cor:mixture}--\ref{cor:anchor-success}. \\
(У3) Почему последовательность шагов приходит к успеху? & Теоремы~\ref{thm:md} и~\ref{thm:md-inexact}: неточный зеркальный спуск; предложение~\ref{prop:ceiling}: потолок представления. \\
(У4) Откуда берётся сигнал и зачем простые навыки? & Лемма~\ref{lem:stationary}, предложения~\ref{prop:barrier}, \ref{prop:dense}, теоремы~\ref{thm:ladder}, \ref{thm:composition}, предложение~\ref{prop:sil}. \\
Ответ целиком & Теорема~\ref{thm:main} и предложение~\ref{prop:necessity}. \\
\bottomrule
\end{tabularx}
\end{center}

Стандартные результаты, на которые опирается глава, явно атрибутированы (раздел~\ref{sec:limits}). Собственный вклад составляют их формулировка для вероятности успеха на конечном горизонте с явными константами, следствия о сохранении навыков и якоре, теорема о лестнице стартов, теорема о композиции с членами переноса инструкции и рассогласования передачи, нижняя граница навыков при сбалансированной самоимитации и сводная теорема.

%======================================================================
\section{Математическая модель}\label{sec:model}

\subsection{Контексты, состояния и история}
Пусть $\cL$ --- конечное множество инструкций, $\Xi$ --- пространство сцен (объекты, их положения, приёмник). Контекст $c=(\ell,\xi)$ распределён по закону $\cD$. Физическое состояние симулятора $s\in S$; наблюдение $o_h$ (изображения камер и проприоцепция) порождается из $s_h$. Решения принимаются в моменты $h=0,\dots,H-1$. Одно решение --- это кусок действий (action chunk), из которого исполняется префикс, после чего политика заново наблюдает сцену. Поэтому $H$ считает \emph{решения}, а не шаги физики; в экспериментах CDPR $H$ порядка $40$--$128$.

Замороженный SmolVLA выдаёт приор действий $p_h$ по наблюдению. Он может быть случайным (в нём собственный шум сэмплирования), но его закон не зависит от обучаемых параметров. Поэтому приор относится к \emph{окружению} как часть канала наблюдения.

Обозначим через $u_h\in\R^m$ латентное (до усечения) действие и через $T(\cdot\mid s,u)$ переходное ядро. В него входят усечение $a=\clip(u,-1,1)$, низкоуровневый контроллер (для CDPR --- обратная кинематика тросов) и физика. \emph{Расширенное состояние}
\[
z_h=(c,\ s_{0:h},\ o_{0:h},\ p_{0:h},\ u_{0:h-1})\in\cZ_h
\]
содержит всю историю. При любой политике, зависящей от истории, процесс $(z_h)$ марковский; $z_H$ есть вся траектория. Все утверждения ниже формулируются для этого процесса. Частичная наблюдаемость учтена тем, что класс политик ограничен (п.~\ref{sec:policy-class}), а не тем, что политика видит $s_h$.

\subsection{Спецификация навыка}
\begin{definition}[Навык и составная задача]\label{def:skill}
\emph{Навык} (задача) $\cT=(\ell,\mu,E,H)$ состоит из инструкции $\ell$, стартового распределения $\mu$ на парах (сцена, начальное состояние), измеримого \emph{события успеха} $E\subseteq\cZ_H$ и горизонта $H$. Вероятность успеха политики $\pi$:
\[
J_\cT(\pi)=\PP^\pi_\mu(z_H\in E).
\]
\emph{Составная задача} $\cT_c=(\ell_c,\mu_c,E_c,H)$ --- задача, чей успех требует последовательного выполнения этапов. \emph{Вехи} $E^{(1)}\supseteq E^{(2)}\supseteq\dots\supseteq E^{(m)}=E_c$ --- вложенные события «этапы $1,\dots,j$ выполнены».
\end{definition}

Примеры из пайплайна: \emph{move to $\langle o\rangle$} (подход с пустым захватом), \emph{pick up $\langle o\rangle$} (стартовое распределение может включать «выровненный захват» над объектом), \emph{put $\langle o\rangle$ into bowl} со стартами «объект уже в захвате» (caught starts), и составная задача с пустым захватом в начале.

\subsection{Класс политик и параметризация}\label{sec:policy-class}
\begin{assumption}[Стационарное представление]\label{ass:features}
Существует фиксированное (не зависящее от обучаемых параметров) измеримое отображение $\varphi_h:\cZ_h\to\R^n$, $x_h=\varphi_h(z_h)$. Оно включает случайную проекцию визуальных токенов, проприоцепцию, приор $p_h$ и инструкцию. Класс допустимых политик:
\[
\Pi_\varphi=\bigl\{\pi=(\pi_h)_{h<H}:\ \pi_h(\cdot\mid z)\ \text{зависит от $z$ только через}\ \varphi_h(z)\bigr\}.
\]
\end{assumption}

Допущение выполняется буквально, когда бэкбон, LoRA и проекция заморожены, как на этапе полной задачи в CDPR. Существенно, что \emph{сравнительная} политика в оценках тоже берётся из $\Pi_\varphi$: никакой привилегированной информации у неё нет.

Реализованная параметризация (CDPR; для WidowX аналогично):
\begin{equation}\label{eq:param}
\mu_\theta(x)=\tanh\bigl(p+s_r\tanh f_\theta(x)\bigr),\qquad
u_h\sim\mathcal N\bigl(\mu_\theta(x_h)+b_h,\ \operatorname{diag}\sigma_\theta^2\bigr),\qquad
a_h=\clip(u_h,-1,1),
\end{equation}
где $f_\theta$ --- остаточная сеть, $s_r$ --- масштаб остатка, $b_h$ --- экзогенное смещение эпизода. Действием в смысле теории считается латентное $u_h$, а усечение отнесено к окружению. Тогда политика имеет полный носитель, и все KL-расстояния ниже конечны. Это согласуется с реализацией, где правдоподобие вычисляется для сохранённого $u_h$ при сохранённом $b_h$.

\subsection{Функции ценности и распределения}
Для политики $\pi$ и задачи $\cT$ положим
\[
V^\pi_h(z)=\PP^\pi(E\mid z_h=z),\quad Q^\pi_h(z,u)=\PP^\pi(E\mid z_h=z,u_h=u),\quad A^\pi_h=Q^\pi_h-V^\pi_h,
\]
так что $V^\pi_H(z)=\one\{z\in E\}$ для всех $\pi$ и $Q,V\in[0,1]$. Через $d^\pi_h$ обозначим закон $z_h$ при $\pi$. Для пары политик $\pi,\pi'$ введём
\[
\alpha_h(z)=\TV\bigl(\pi'_h(\cdot\mid z),\pi_h(\cdot\mid z)\bigr),\qquad
\ab_h=\E_{z\sim d^\pi_h}\alpha_h(z),\qquad
\bar A_h(z)=\E_{u\sim\pi'_h(\cdot\mid z)}A^\pi_h(z,u),
\]
где $\TV(P,Q)=\sup_B|P(B)-Q(B)|$. Ключевое обстоятельство: \emph{$\ab_h$ вычисляется по состояниям старой политики}, то есть по тем данным, на которых обновление и строится.

\subsection{Алгоритм (идеализация GRPO)}\label{sec:algo}
На обновлении $k$ выбираются контексты. Из каждого стартового состояния запускается группа из $G$ независимых эпизодов с исходами $Y_i\in\{0,1\}$. Групповое преимущество
\begin{equation}\label{eq:grpo-adv}
A_i=\frac{Y_i-\bar Y}{\max(\hat\sigma,\epsilon)},\qquad \bar Y=\tfrac1G\textstyle\sum_jY_j,\quad \hat\sigma^2=\tfrac1G\sum_j(Y_j-\bar Y)^2,
\end{equation}
и усечённая суррогатная функция
\begin{equation}\label{eq:surrogate}
\cL^{\rm clip}(\theta)=\frac1N\sum_{i}\sum_{h}\min\Bigl(r_{ih}(\theta)A_i,\ \clip\bigl(r_{ih}(\theta),1-\varepsilon_-,1+\varepsilon_+\bigr)A_i\Bigr),\quad r_{ih}=\frac{\pi_\theta(u_{ih}\mid x_{ih})}{\pi_{\theta_k}(u_{ih}\mid x_{ih})}.
\end{equation}
Её максимизируют несколькими эпохами по мини-батчам, с энтропийным бонусом и, возможно, с якорем к эталонной политике. В конфигурации WidowX: $G=8$, $\varepsilon_-=0{,}2$, $\varepsilon_+=0{,}28$, 4 эпохи, коэффициент энтропии $2\cdot10^{-4}$, фильтр групп по стандартному отклонению $0{,}05$ (динамическое сэмплирование в духе DAPO~\cite{dapo}). Для CDPR кредит дополнительно распределяется по вехам~\cite{notes}. Оптимизатор Adam, усечение нормы градиента и разделение на мини-батчи в теории не моделируются. Их влияние входит в ошибку шага (теорема~\ref{thm:md-inexact}).

%======================================================================
\section{Вспомогательные утверждения}\label{sec:lemmas}

\begin{lemma}[Разность производительностей, конечный горизонт]\label{lem:pdl}
Для любых политик $\pi,\pi'$ и задачи $\cT$
\[
J(\pi')-J(\pi)=\sum_{h=0}^{H-1}\E_{z\sim d^{\pi'}_h}\bar A_h(z)
=\sum_{h=0}^{H-1}\E_{z\sim d^{\pi'}_h}\ip{\pi'_h(\cdot\mid z)-\pi_h(\cdot\mid z)}{Q^\pi_h(z,\cdot)}.
\]
\end{lemma}
\begin{proof}
Так как $V^\pi_H=\one_E$ и $d_0$ не зависит от политики, имеем $J(\pi')=\E^{\pi'}V^\pi_H(z_H)$ и $J(\pi)=\E^{\pi'}V^\pi_0(z_0)$. Телескопируя,
$J(\pi')-J(\pi)=\sum_h\E^{\pi'}[V^\pi_{h+1}(z_{h+1})-V^\pi_h(z_h)]$. Переходное ядро не зависит от политики, поэтому
$\E[V^\pi_{h+1}(z_{h+1})\mid z_h,u_h]=Q^\pi_h(z_h,u_h)$, и слагаемое равно $\E^{\pi'}[Q^\pi_h(z_h,u_h)-V^\pi_h(z_h)]=\E_{d^{\pi'}_h}\bar A_h$. Второе равенство следует из $\E_{u\sim\pi_h}A^\pi_h(z,u)=0$.
\end{proof}

\begin{lemma}[Сдвиг распределения состояний]\label{lem:shift}
Для $h=0,\dots,H$
\[
\TV\bigl(d^{\pi'}_h,d^\pi_h\bigr)\le\sum_{t<h}\ab_t .
\]
В частности, для законов траекторий $\TV(\PP^{\pi'},\PP^\pi)\le\sum_{t=0}^{H-1}\ab_t$.
\end{lemma}
\begin{proof}
Пусть $K^\pi_t(z,\cdot)=\int\pi_t(du\mid z)\,T_t(\cdot\mid z,u)$ --- ядро перехода $z_t\to z_{t+1}$. Тогда
$d^{\pi'}_{t+1}-d^\pi_{t+1}=(d^{\pi'}_t-d^\pi_t)K^{\pi'}_t+d^\pi_t(K^{\pi'}_t-K^\pi_t)$.
Марковское ядро не увеличивает расстояние по вариации, поэтому TV-норма первого слагаемого не превосходит $\TV(d^{\pi'}_t,d^\pi_t)$. Для второго по неравенству обработки данных (одно и то же ядро $T_t(\cdot\mid z,\cdot)$ применяется к двум законам действий)
$\|d^\pi_t(K^{\pi'}_t-K^\pi_t)\|_{\TV}\le\int d^\pi_t(dz)\TV(\pi'_t(\cdot\mid z),\pi_t(\cdot\mid z))=\ab_t$.
Индукция по $t$ с $\TV(d^{\pi'}_0,d^\pi_0)=0$. Поскольку $z_H$ --- вся траектория, $d_H$ есть закон траектории.
\end{proof}

\begin{lemma}[Пинскер в среднем]\label{lem:pinsker}
$\ab_h\le\sqrt{\tfrac12\E_{z\sim d^\pi_h}\KL(\pi_h(\cdot\mid z)\,\|\,\pi'_h(\cdot\mid z))}$; то же верно для обратного направления KL.
\end{lemma}
\begin{proof}
Неравенство Пинскера $\TV(P,Q)\le\sqrt{\KL(P\|Q)/2}$ и неравенство Йенсена для вогнутой функции $\sqrt{\cdot}$.
\end{proof}

Лемма~\ref{lem:pinsker} связывает теорию с журналами обучения. Средний выборочный KL на обновлении (\texttt{approx\_kl\_mean}) оценивает $\E_{d^\pi}\KL$ на состояниях старой политики и тем самым даёт верхнюю оценку $\ab_h$.

%======================================================================
\section{Что оптимизирует GRPO при бинарной награде}\label{sec:objective}

При $R=\one_E$ целевая функция RL совпадает с метрикой задачи: $J(\pi)=\E_c\PP^\pi(E\mid c)$. Суррогатной награды нет. Покажем, что GRPO в среднем поднимается по монотонному преобразованию этой вероятности.

\begin{theorem}[Эффективная цель GRPO]\label{thm:arcsin}
Зафиксируем контекст $c$ и положим $p_c(\theta)=\PP^{\pi_\theta}(E\mid c)\in(0,1)$. Пусть функция правдоподобия траектории дифференцируема, а скор $Z=\sum_h\nabla_\theta\log\pi_\theta(u_h\mid x_h)$ имеет конечный второй момент (для гауссовской политики~\eqref{eq:param} с гладким средним и отделённой от нуля дисперсией на конечном горизонте это выполнено). Рассмотрим направление обновления в точке поведения $\theta=\theta_k$, где все $r_{ih}=1$ и усечение неактивно: $g=\frac1G\sum_iA_iZ_i$.
\begin{enumerate}[label=(\alph*)]
\item При нормировке по генеральной совокупности, $A=(Y-p)/\sqrt{p(1-p)}$,
\[
\E g=\frac{\nabla_\theta p_c}{\sqrt{p_c(1-p_c)}}=\nabla_\theta\bigl[2\arcsin\sqrt{p_c(\theta)}\bigr].
\]
\item При групповой нормировке~\eqref{eq:grpo-adv} $\E g=c_{G,\epsilon}(p_c)\nabla_\theta p_c$, где
\[
c_{G,\epsilon}(p)=\frac{1}{p(1-p)}\,\E\Bigl[\frac{Q(1-Q)}{\max(\sqrt{Q(1-Q)},\epsilon)}\Bigr]>0,\qquad Q=K/G,\quad K\sim\mathrm{Bin}(G,p);
\]
при $G\to\infty$ имеем $c_{G,\epsilon}(p)\to1/\sqrt{p(1-p)}$.
\item Для контекстов $c\sim\cD$ с общими параметрами $\E g=\nabla_\theta\E_cF_G(p_c(\theta))$, где $F_G'=c_{G,\epsilon}$, и $F_G$ строго возрастает на $(0,1)$.
\end{enumerate}
\end{theorem}
\begin{proof}
Производная $\min(rA,\clip(r,1-\varepsilon_-,1+\varepsilon_+)A)$ по $\theta$ при $r=1$ равна $A\nabla_\theta r=A\nabla_\theta\log\pi_\theta$, так как $1$ --- внутренняя точка интервала усечения; суммирование по $h$ даёт $AZ$.
(а) По тождеству скор-функции $\E[Z]=0$ и $\E[YZ]=\nabla\E[Y]=\nabla p$, поэтому $\E[(Y-p)Z]=\nabla p$. Остаётся заметить, что $\frac{d}{dp}2\arcsin\sqrt p=1/\sqrt{p(1-p)}$.
(б) Условно на метки $Y_{1:G}$ траектории независимы, поэтому $\E[Z_i\mid Y_{1:G}]=\E[Z_i\mid Y_i]$; при этом $\E[Z\mid Y=1]=\nabla p/p$ и $\E[Z\mid Y=0]=-\nabla p/(1-p)$. Подстановка даёт
$\E[g\mid Y_{1:G}]=\frac{Q(1-Q)}{d(Q)}\bigl(\frac1p+\frac1{1-p}\bigr)\nabla p$ с $d(Q)=\max(\sqrt{Q(1-Q)},\epsilon)$; осталось усреднить по $K$. Предел следует из $Q\to p$ п.н. и ограниченности подынтегрального выражения. Формула и её связь с~(а) получены в~\cite{notes}; см. также~\cite{mroueh}.
(в) Линейность математического ожидания; закон выбора контекстов от $\theta$ не зависит.
\end{proof}

\begin{corollary}[Свойства эффективной цели]\label{cor:arcsin}
\begin{enumerate}[label=(\roman*)]
\item Контекст с $p_c\in\{0,1\}$ не даёт вклада. В реализации вырожденная группа даёт $A_i\equiv0$ точно, а не только в среднем.
\item Относительно обычного градиента политики вес контекста равен $1/\sqrt{p(1-p)}$: трудные контексты ($p\to0$) усиливаются как $p^{-1/2}$. Функция $2\arcsin\sqrt p$ вогнута на $[0,\tfrac12]$ и выпукла на $[\tfrac12,1]$.
\item \emph{Парето-оптимальность.} Если $\theta^\ast$ --- (локальный) максимум $\E_cF_G(p_c(\theta))$, то вектор $(p_c(\theta^\ast))_c$ (локально) Парето-оптимален: не существует $\theta'$ (в окрестности) с $p_c(\theta')\ge p_c(\theta^\ast)$ для $\cD$-п.в.\ $c$ и строгим неравенством на множестве положительной меры.
\end{enumerate}
\end{corollary}
\begin{proof}
(i) и (ii) следуют из теоремы~\ref{thm:arcsin}. (iii) Если бы такое $\theta'$ существовало, строгое возрастание $F_G$ дало бы $\E_cF_G(p_c(\theta'))>\E_cF_G(p_c(\theta^\ast))$, что противоречит максимальности.
\end{proof}

Таким образом, GRPO оптимизирует не арифметическое среднее успеха, а другую скаляризацию того же вектора вероятностей успеха по контекстам. Её оптимумы Парето-оптимальны, а если в классе есть политика, решающая все контексты, обе скаляризации достигают максимума на ней.

\begin{remark}[Константы реализации]\label{rem:constants}
Для бинарной группы размера $G$ максимум $|A_i|$ равен $\sqrt{G-1}$ (один успех из $G$); при $G=8$ это $2{,}65$. Поэтому порог усечения преимуществ $6{,}0$ при бинарной награде не срабатывает. Минимальное ненулевое групповое стандартное отклонение равно $\sqrt{(G-1)}/G\approx0{,}33>0{,}05$, так что фильтр по стандартному отклонению удаляет ровно вырожденные группы.
\end{remark}

\begin{remark}[Вехи]
Для вложенных вех $E^{(1)}\supseteq\dots\supseteq E^{(m)}$ с независимо нормированными преимуществами и весами $w_j\ge0$ то же рассуждение даёт среднее направление $\sum_jw_jc_{G,\epsilon}(P(E^{(j)}))\nabla P(E^{(j)})$, а по цепному правилу $J=\prod_j P(E^{(j)}\mid E^{(j-1)})$. Анализ правила защиты достигнутых вех ($\kappa=0$) дан в~\cite[предложение~C]{notes}.
\end{remark}

%======================================================================
\section{Устойчивость: доверительная область и сохранение навыков}\label{sec:stability}

\begin{theorem}[Улучшение в доверительной области для вероятности успеха]\label{thm:improve}
Для любых политик $\pi,\pi'$ положим $L_\pi(\pi')=\sum_h\E_{d^\pi_h}\bar A_h$ и $\varepsilon=\sup_{h,z}|\bar A_h(z)|$. Тогда
\begin{equation}\label{eq:improve}
J(\pi')\ \ge\ J(\pi)+L_\pi(\pi')-2\varepsilon\sum_{h=1}^{H-1}\sum_{t<h}\ab_t
\ \ge\ J(\pi)+L_\pi(\pi')-H(H-1)\,\varepsilon\max_t\ab_t .
\end{equation}
Кроме того, $\varepsilon\le\alpha_\infty:=\sup_{h,z}\alpha_h(z)$, откуда
$J(\pi')\ge J(\pi)+L_\pi(\pi')-H(H-1)\alpha_\infty^2\ge J(\pi)+L_\pi(\pi')-\tfrac12H(H-1)\delta$,
где $\delta=\sup_{h,z}\KL(\pi_h(\cdot\mid z)\|\pi'_h(\cdot\mid z))$.
\end{theorem}
\begin{proof}
По лемме~\ref{lem:pdl} $J(\pi')-J(\pi)-L_\pi(\pi')=\sum_h\int\bar A_h\,d(d^{\pi'}_h-d^\pi_h)$. Для функции со значениями в $[-\varepsilon,\varepsilon]$ имеем $|\int f\,d(\nu-\nu')|\le2\varepsilon\TV(\nu,\nu')$, а по лемме~\ref{lem:shift} $\TV(d^{\pi'}_h,d^\pi_h)\le\sum_{t<h}\ab_t$. Слагаемое с $h=0$ равно нулю, и $\sum_{h=1}^{H-1}h=H(H-1)/2$. Далее, $\bar A_h(z)=\E_{\pi'_h}Q^\pi_h(z,\cdot)-\E_{\pi_h}Q^\pi_h(z,\cdot)$, а $Q\in[0,1]$, поэтому $|\bar A_h(z)|\le\alpha_h(z)$. Последнее неравенство --- неравенство Пинскера.
\end{proof}

\begin{corollary}[Монотонность шага миноризации--максимизации]\label{cor:mm}
Положим $M_\pi(\pi')=J(\pi)+L_\pi(\pi')-\tfrac12H(H-1)\delta(\pi,\pi')$. Тогда $J(\pi')\ge M_\pi(\pi')$ и $M_\pi(\pi)=J(\pi)$. Следовательно, любое обновление с $M_{\pi_k}(\pi_{k+1})\ge M_{\pi_k}(\pi_k)$ не уменьшает вероятность успеха: $J(\pi_{k+1})\ge J(\pi_k)$.
\end{corollary}

Это аналог рассуждения TRPO~\cite{trpo} для вероятности успеха на конечном горизонте; множитель $H(H-1)$ при $\alpha_\infty^2$ играет роль множителя $4\varepsilon\gamma/(1-\gamma)^2$ из дисконтированного случая (там $\varepsilon=\max|A|$), и при эффективном горизонте $1/(1-\gamma)\approx H$ оба растут как $H^2$. CPO-подобная форма с ожидаемой вариацией приведена в~\cite{cpo}.

\begin{remark}[Как это соотносится с усечением и что из этого следует]\label{rem:clip}
\begin{enumerate}[label=(\roman*)]
\item Если бы для всех действий отношение $r=\pi'/\pi$ лежало в $[1-\varepsilon_-,1+\varepsilon_+]$, то $\alpha_h(z)=\tfrac12\E_\pi|r-1|\le\tfrac12\max(\varepsilon_-,\varepsilon_+)=0{,}14$. Усечение PPO этого не гарантирует: оно лишь убирает стимул выходить за полосу на \emph{выборочных} действиях~\cite{engstrom,truly-ppo}. Поэтому условия теоремы --- свойство \emph{реализованного} шага, и его надо контролировать по журналу (\texttt{approx\_kl}, доля усечений, отбраковка шагов с большим KL).
\item Допустимый шаг убывает как $H^{-2}$: улучшение гарантировано при $\delta<2L_\pi(\pi')/(H(H-1))$. Отсюда теоретический довод в пользу кусков действий (уменьшают $H$ как число решений) и малых шагов на длинных задачах.
\item Константы худшего случая: при $H=40$ и $\delta=0{,}009$ штраф $\tfrac12H(H-1)\delta\approx7$ превосходит единицу, и оценка численно бессодержательна. Её значение структурное: она объясняет \emph{почему} и \emph{как} размер шага должен зависеть от горизонта, но не сертифицирует конкретное обновление.
\end{enumerate}
\end{remark}

Следующая теорема использует только средние по состояниям старой политики. Поэтому она проверяема по данным, и численная проверка показывает, что она почти точна (приложение~\ref{app:numerics}).

\begin{theorem}[Сохранение навыка]\label{thm:retention}
Для любой задачи $\cT$ и любых $\pi,\pi'$
\[
\bigl|J_\cT(\pi')-J_\cT(\pi)\bigr|\le\TV\bigl(\PP^{\pi'}_\cT,\PP^\pi_\cT\bigr)\le\Delta_\cT(\pi,\pi'):=\sum_{h=0}^{H-1}\E_{z\sim d^{\pi,\cT}_h}\alpha_h(z)
\le\sum_{h=0}^{H-1}\sqrt{\tfrac12\E_{d^{\pi,\cT}_h}\KL\bigl(\pi_h\|\pi'_h\bigr)} .
\]
\end{theorem}
\begin{proof}
$E$ --- событие, поэтому $|\PP^{\pi'}(E)-\PP^\pi(E)|\le\TV$. Далее применяются лемма~\ref{lem:shift} при $h=H$ и лемма~\ref{lem:pinsker}.
\end{proof}

Принципиально, что ожидание берётся по состояниям \emph{данной} задачи $\cT$ при старой политике. Обновление контролирует изменение политики только там, где оно собирает данные.

\begin{corollary}[Смесь задач и цена исключения]\label{cor:mixture}
Пусть обучение ведётся на смеси задач с весами $w_\cT\ge0$, $\sum w_\cT=1$, и реализованный шаг удовлетворяет $\sum_\cT w_\cT\Delta_\cT(\pi,\pi')\le\bar\Delta$. Тогда для каждой задачи с $w_\cT>0$
\[
|J_\cT(\pi')-J_\cT(\pi)|\le\bar\Delta/w_\cT .
\]
Для задачи с $w_\cT=0$ никакой оценки не следует: её изменение определяется обобщением общей сети и не контролируется доверительной областью.
\end{corollary}
\begin{proof}
Все слагаемые неотрицательны, поэтому $w_\cT\Delta_\cT\le\bar\Delta$; далее применяется теорема~\ref{thm:retention}.
\end{proof}

\begin{corollary}[Накопление дрейфа]\label{cor:drift}
После $K$ обновлений $|J_\cT(\pi_K)-J_\cT(\pi_0)|\le\sum_{k<K}\Delta_\cT(\pi_k,\pi_{k+1})$. Оценка растёт линейно по $K$, то есть доверительная область на каждом шаге \emph{не} ограничивает суммарную деградацию. Более того, теорема~\ref{thm:improve} для смесевой цели $J_w=\sum w_\cT J_\cT$ допускает «обмен» навыков: $J_w$ растёт, тогда как отдельная $J_\cT$ падает, причём на каждом шаге не более чем на $\bar\Delta/w_\cT$.
\end{corollary}

\begin{corollary}[Якорь к эталонной политике]\label{cor:anchor}
Пусть $\pi_{\rm ref}$ --- замороженная эталонная политика и
$D_{\rm ref}(\pi)=\sum_h\E_{z\sim d^{\pi_{\rm ref},\cT}_h}\KL(\pi_{{\rm ref},h}(\cdot\mid z)\|\pi_h(\cdot\mid z))$. Тогда для \emph{любой} итерации, независимо от числа шагов,
\[
|J_\cT(\pi)-J_\cT(\pi_{\rm ref})|\le\sqrt{H\,D_{\rm ref}(\pi)/2}.
\]
Для гауссовских политик с общей ковариацией $\KL=\sum_j(\mu_j-\mu_{{\rm ref},j})^2/(2\sigma_j^2)$; именно этот член реализован в якоре (\texttt{reference\_anchor.py}).
\end{corollary}
\begin{proof}
Применим теорему~\ref{thm:retention}, взяв $\pi_{\rm ref}$ в роли старой политики: ожидания берутся по $d^{\pi_{\rm ref}}$. Затем применим Пинскера и неравенство Коши--Буняковского: $\sum_h\sqrt{x_h/2}\le\sqrt{H\sum_hx_h/2}$.
\end{proof}

В реализации банк якоря содержит не все состояния эталона, а состояния его \emph{успешных} эпизодов. Для этого случая есть односторонняя версия.

\begin{corollary}[Сохранение по банку успехов]\label{cor:anchor-success}
Пусть $q^+$ --- закон траекторий $\pi_{\rm ref}$ при условии успеха и
$\Lambda(\pi)=\E_{q^+}\sum_h\log\frac{\pi_{{\rm ref},h}(u_h\mid z_h)}{\pi_h(u_h\mid z_h)}$. Тогда
\[
J_\cT(\pi)\ \ge\ J_\cT(\pi_{\rm ref})\,e^{-\Lambda(\pi)}\ \ge\ J_\cT(\pi_{\rm ref})\bigl(1-\Lambda(\pi)\bigr).
\]
Для гауссовских политик с общей диагональной ковариацией
$\Lambda=\sum_h\E_{q^+}\sum_j\Bigl[\frac{(\mu_j-\mu_{{\rm ref},j})^2}{2\sigma_j^2}+\frac{(\mu_{{\rm ref},j}-\mu_j)(u_j-\mu_{{\rm ref},j})}{\sigma_j^2}\Bigr]$.
Первый член есть штраф якоря на банке успехов. Второй имеет первый порядок по $\|\mu-\mu_{\rm ref}\|$ и исчезает, если на состояниях банка успех не смещает латентное действие относительно среднего эталона.
\end{corollary}
\begin{proof}
$J(\pi)\ge\int_E d\PP^\pi=J(\pi_{\rm ref})\E_{q^+}[d\PP^\pi/d\PP^{\pi_{\rm ref}}]\ge J(\pi_{\rm ref})\exp(\E_{q^+}\log(d\PP^\pi/d\PP^{\pi_{\rm ref}}))$ по неравенству Йенсена. Множители динамики в отношении правдоподобий сокращаются. Формула для гауссовского случая получается раскрытием квадратов.
\end{proof}

%======================================================================
\section{Сходимость: консервативные шаги как неточный зеркальный спуск}\label{sec:convergence}

Теоремы раздела~\ref{sec:stability} гарантируют, что шаг не вредит. Теперь покажем, что последовательность таких шагов \emph{приходит} к успешной политике.

\begin{definition}[KL-проксимальный шаг]\label{def:md}
Для оценки $\hat Q_k$ шаг
\[
\pi_{k+1,h}(\cdot\mid z)=\argmax_q\Bigl\{\E_{u\sim q}\hat Q_{k,h}(z,u)-\tfrac1\eta\KL\bigl(q\,\|\,\pi_{k,h}(\cdot\mid z)\bigr)\Bigr\}
=\frac{\pi_{k,h}(\cdot\mid z)\,e^{\eta\hat Q_{k,h}(z,\cdot)}}{\E_{\pi_{k,h}}e^{\eta\hat Q_{k,h}(z,\cdot)}}
\]
(второе равенство --- вариационный принцип Гиббса) называется \emph{KL-проксимальным}, или шагом зеркального спуска в пространстве политик.
\end{definition}

Суррогат PPO с KL-штрафом в популяционной форме совпадает с целевой функцией этого шага при $\eta=1/\beta$. Усечённый суррогат~\eqref{eq:surrogate} есть его приближение первого порядка, где полоса усечения играет роль $1/\eta$~\cite{mdpo,lan,shani,neural-ppo}. GRPO подставляет в качестве $\hat Q$ монте-карловские исходы группы: $\E[Y\mid z_h,u_h]=Q^\pi_h(z_h,u_h)$.

\begin{theorem}[Сходимость к любому сравнимому поведению]\label{thm:md}
Пусть $\hat Q_k=Q^{\pi_k}$ (точные оценки) и $\pit$ --- произвольная политика с
\[
D(\pit\|\pi_0):=\sum_{h=0}^{H-1}\E_{z\sim d^{\pit}_h}\KL\bigl(\pit_h(\cdot\mid z)\,\|\,\pi_{0,h}(\cdot\mid z)\bigr)<\infty .
\]
Тогда для любых $\eta>0$ и $K\ge1$
\[
\frac1K\sum_{k<K}\bigl[J(\pit)-J(\pi_k)\bigr]\le\frac{D(\pit\|\pi_0)}{\eta K}+\frac{H\eta}8,
\]
и при $\eta=\sqrt{8D/(HK)}$
\[
\min_{k<K}\bigl[J(\pit)-J(\pi_k)\bigr]\le\sqrt{\frac{H\,D(\pit\|\pi_0)}{2K}} .
\]
\end{theorem}
\begin{proof}
Зафиксируем $(h,z)$ и обозначим $Q=Q^{\pi_k}_h(z,\cdot)\in[0,1]$, $Z_k=\E_{\pi_k}e^{\eta Q}$. Из определения шага для любого закона $q$
\[
\KL(q\|\pi_{k+1})-\KL(q\|\pi_k)=\E_q\log\frac{\pi_k}{\pi_{k+1}}=-\eta\E_qQ+\log Z_k .
\]
По лемме Хёфдинга $\log Z_k\le\eta\E_{\pi_k}Q+\eta^2/8$. При $q=\pit_h(\cdot\mid z)$ получаем
\[
\ip{\pit_h-\pi_{k,h}}{Q^{\pi_k}_h}(z)\le\frac{\KL(\pit\|\pi_k)(z)-\KL(\pit\|\pi_{k+1})(z)}{\eta}+\frac\eta8 .
\]
По лемме~\ref{lem:pdl} с $\pi'=\pit$ имеем $J(\pit)-J(\pi_k)=\sum_h\E_{d^{\pit}_h}\ip{\pit_h-\pi_{k,h}}{Q^{\pi_k}_h}$. Ключевое обстоятельство: мера $d^{\pit}_h$ \emph{не зависит от $k$}. Поэтому суммирование по $k<K$ телескопирует KL-члены в каждой точке $z$:
\[
\sum_{k<K}\bigl[J(\pit)-J(\pi_k)\bigr]\le\frac1\eta\sum_h\E_{d^{\pit}_h}\bigl[\KL(\pit\|\pi_0)-\KL(\pit\|\pi_K)\bigr]+\frac{KH\eta}8\le\frac{D}{\eta}+\frac{KH\eta}8 .
\]
Минимизация по $\eta$ завершает доказательство.
\end{proof}

Теорема не требует, чтобы сравнительная политика была оптимальной: достаточно любого успешного поведения на конечном KL-расстоянии от начальной политики. Это и есть ответ на подвопрос~(У3) в идеализированной форме. Консервативные шаги не мешают сходимости, их размер ($\eta$) лишь задаёт скорость. Вся «трудность задачи» сосредоточена в одном числе $D$.

\begin{corollary}[Гауссовские политики: роль приора и начального шума]\label{cor:gauss}
Пусть $\pi_{0,h}(\cdot\mid z)=\mathcal N(\mu_0(x),\sigma_0^2I_m)$, $\pit_h(\cdot\mid z)=\mathcal N(\tilde\mu(x),\tilde\sigma^2I_m)$. Тогда
\begin{multline*}
D(\pit\|\pi_0)=\sum_h\E_{d^{\pit}_h}\frac{\|\tilde\mu(x)-\mu_0(x)\|^2}{2\sigma_0^2}\\
+Hm\Bigl[\log\frac{\sigma_0}{\tilde\sigma}+\frac{\tilde\sigma^2}{2\sigma_0^2}-\frac12\Bigr],
\end{multline*}
и для средней ошибки $\epsilon$ достаточно $K\ge HD/(2\epsilon^2)$ шагов. Отсюда:
(а) скорость определяется рассогласованием среднего начальной политики с успешным поведением \emph{вдоль траекторий успешного поведения}, в единицах начального шума;
(б) требуемая точность $\tilde\sigma\ll\sigma_0$ обходится лишь логарифмически;
(в) нулевая инициализация остатка ставит $\mu_0=\tanh(p)$, то есть приор. Если приор --- постоянное смещение, далёкое от нужных действий (как на CDPR), $D$ велико: остаток должен его компенсировать и работает в насыщении $\tanh$.
\end{corollary}
\begin{proof}
Формула KL двух изотропных гауссиан и теорема~\ref{thm:md}.
\end{proof}

Например, при $\sigma_0=e^{-1{,}2}\approx0{,}30$ (начальное значение $\log\sigma$ в реализации) рассогласование $0{,}1$ по одной координате стоит $0{,}055$ нат на шаг, а $0{,}5$ --- уже $1{,}38$ нат. Это в 25 раз больше, и при той же точности требуется в 25 раз больше обновлений. Ограниченность~\eqref{eq:param} добавляет условие представимости: среднее сравнительной политики должно лежать в $[\tanh(p-s_r),\tanh(p+s_r)]$~\cite[\S6]{notes}.

Реальный шаг не является точным шагом Гиббса: $Q$ оценивается по $G$ исходам, шаг проецируется на класс сети над замороженными признаками, а усечение лишь приближает проксимальный шаг. Все эти отклонения учитывает следующая теорема.

\begin{theorem}[Неточный шаг: оценивание, представление и покрытие]\label{thm:md-inexact}
Пусть $\pi_k$ --- произвольная последовательность политик. Определим \emph{невязку шага}
\[
e_{k,h}(z)=\ip{\pit_h-\pi_{k,h}}{Q^{\pi_k}_h}(z)-\frac{\KL(\pit\|\pi_k)(z)-\KL(\pit\|\pi_{k+1})(z)}{\eta}-\frac\eta8
\]
(для точного шага Гиббса $e_{k,h}\le0$). Пусть $\nu_h$ --- распределение состояний, на которых собираются данные шага $h$ (по политике и из рестартов), и $C_h=\E_{\nu_h}\bigl[(dd^{\pit}_h/d\nu_h)^2\bigr]=1+\chi^2(d^{\pit}_h\|\nu_h)$. Тогда
\[
\frac1K\sum_{k<K}\bigl[J(\pit)-J(\pi_k)\bigr]\le\frac{D(\pit\|\pi_0)}{\eta K}+\frac{H\eta}8+\frac1K\sum_{k<K}\sum_h\sqrt{C_h}\,\bigl\|e^+_{k,h}\bigr\|_{L^2(\nu_h)} .
\]
\end{theorem}
\begin{proof}
Как в теореме~\ref{thm:md}, но с невязкой. По неравенству Коши--Буняковского
\[
\E_{d^{\pit}_h}e_{k,h}\le\E_{\nu_h}\Bigl[\frac{dd^{\pit}_h}{d\nu_h}\,e^+_{k,h}\Bigr]\le\sqrt{C_h}\,\|e^+_{k,h}\|_{L^2(\nu_h)} .
\]
\end{proof}

Три члена оценки отвечают трём механизмам. \emph{Оптимизационный} член $D/(\eta K)+H\eta/8$ отражает то, что консервативный шаг сходится. \emph{Невязка} $e$ складывается из ошибки оценивания $Q$ по группе (на вырожденной группе оценка преимущества равна нулю, и невязка равна полному истинному преимуществу), ошибки представления (шаг Гиббса не выразим на признаках $\varphi$) и ошибки приближения проксимального шага усечённым суррогатом. Член \emph{покрытия} $C_h$ показывает, насколько состояния, на которых учимся, покрывают состояния успешного поведения. Если $\nu_h$ не имеет массы там, где её имеет $d^{\pit}_h$, то $C_h=\infty$ и гарантии нет: алгоритм не знает, что делать в состояниях, которых не видел. Такова формальная роль стартовых распределений (объект в захвате, выровненный захват, лестница расстояний) и рестартов из промежуточных состояний~\cite{kakade-langford,agarwal}.

\begin{remark}[Реактивные политики]\label{rem:reactive}
Расширенное состояние содержит время и историю, поэтому рестарт «объект в захвате» в момент $0$ формально не совпадает с состоянием составной задачи в момент $h$. Если политика реактивна (зависит от текущего наблюдения, приора и инструкции; внутреннюю память можно включить в $s$), а событие успеха определяется будущим участком траектории при достаточном оставшемся горизонте, все конструкции переносятся на пространство (физическое состояние, инструкция). Тогда покрытие --- свойство распределения физических состояний, и рестарты непосредственно уменьшают $C_h$.
\end{remark}

\begin{proposition}[Потолок представления]\label{prop:ceiling}
Пусть на фиксированном решающем шаге $h^\ast$ успех требует $\|g(\xi)-u_{h^\ast}\|\le r$, где $g(\xi)$ --- координаты захвата в сцене (с точностью до известного отображения), и пусть $\sup_a\PP(\|g(\xi)-a\|\le r\mid x_{h^\ast})\le\beta$ почти наверное. Тогда $J(\pi)\le\beta$ для любой $\pi\in\Pi_\varphi$, каким бы алгоритмом она ни была получена.
\end{proposition}
\begin{proof}
При $\pi\in\Pi_\varphi$ действие $u_{h^\ast}$ разыгрывается из $\pi(\cdot\mid x_{h^\ast})$ с независимой случайностью, поэтому $u_{h^\ast}\perp g(\xi)\mid x_{h^\ast}$. Следовательно, $\PP(\|g-u\|\le r\mid x)=\int\pi(du\mid x)\PP(\|g-u\|\le r\mid x)\le\beta$; усредняем по $x$.
\end{proof}

Следовательно, теорема~\ref{thm:md} приводит RL к $J^\ast_\varphi=\sup_{\Pi_\varphi}J$, а не к $\sup J$ по всем политикам. Разрыв $\sup J-J^\ast_\varphi$ определяется \emph{представлением}, а не алгоритмом RL, и устраняется только улучшением признаков (адаптация зрения, камера на запястье), но не увеличением бюджета RL. Важна точность признаков в момент решения: в замкнутом контуре она растёт по мере приближения к объекту.

\begin{remark}[Случайная проекция признаков]\label{rem:jl}
Визуальные токены проецируются фиксированной гауссовской матрицей в $\R^{512}$. По лемме Джонсона--Линденштраусса~\cite{jl,dasgupta} для любого конечного множества из $N$ векторов попарные расстояния сохраняются с множителем $1\pm\epsilon$ с вероятностью не меньше $1-\delta$ при $512\gtrsim\epsilon^{-2}\log(N/\delta)$. Поэтому липшицева политика на исходных токенах, ограниченная на это множество, задаётся $L/(1-\epsilon)$-липшицевой функцией на проекциях (по теореме Киршбрауна). Потолок $J^\ast_\varphi$ определяется в основном самими токенами VLM, а проекция ухудшает его не более чем через искажение $\epsilon$.
\end{remark}

\paragraph{Почему обучается только исполнительная часть.} В рамках изложенной теории у этого решения три причины.
(1)~При замороженном бэкбоне класс $\Pi_\varphi$ фиксирован, поэтому сравнительная политика в теоремах~\ref{thm:md} и~\ref{thm:md-inexact} не меняется со временем. При обучении бэкбона класс сдвигается, и анализ становится анализом движущейся цели.
(2)~Входы $(x,p)$ стационарны, поэтому банки самоимитации и якоря (следствия~\ref{cor:anchor}, \ref{cor:anchor-success}, предложение~\ref{prop:sil}) остаются корректными на протяжении всего запуска.
(3)~Геометрия доверительной области задана в пространстве \emph{действий} (KL между политиками) и от параметризации не зависит. Голова --- самый дешёвый способ перемещаться в этой геометрии, а представление остаётся отдельной, явно измеримой переменной (предложение~\ref{prop:ceiling}).

%======================================================================
\section{Сигнал при разреженной награде и роль простых навыков}\label{sec:signal}

\begin{lemma}[Неподвижность до первого успеха]\label{lem:stationary}
Если на обновлении все группы вырождены (все $Y_i$ в группе равны), то $A_i\equiv0$, и вклад градиента политики в обновление равен нулю. Параметры меняются только регуляризаторами (энтропия, якорь), направление которых не зависит от сигнала успеха. При независимых контекстах число групп до первой смешанной группы имеет математическое ожидание $1/\E_c\Pmix(p_c,G)$, где $\Pmix(p,G)=1-p^G-(1-p)^G\le Gp$. Следовательно, ожидаемое число \emph{эпизодов} до первого ненулевого шага не меньше $1/\E_cp_c$.
\end{lemma}
\begin{proof}
При $Y_1=\dots=Y_G$ числитель в~\eqref{eq:grpo-adv} равен нулю. Число групп до первой смешанной имеет геометрическое распределение. Неравенство $1-(1-p)^G\le Gp$ есть неравенство Бернулли.
\end{proof}

\begin{proposition}[Экспоненциальный барьер составной задачи]\label{prop:barrier}
Пусть для вех составной задачи $\PP^{\pi_0}(E^{(j)}\mid E^{(j-1)})\le q<1$ при всех $j=1,\dots,m$. Тогда $J_c(\pi_0)\le q^m$, и ожидаемое число эпизодов до первого ненулевого шага GRPO на составной задаче не меньше $q^{-m}$.
\end{proposition}
\begin{proof}
Цепное правило $J_c=\prod_j\PP(E^{(j)}\mid E^{(j-1)})$ и лемма~\ref{lem:stationary}.
\end{proof}

Модельная иллюстрация. Пусть при исследовательском шуме конечное положение захвата примерно равномерно в шаре радиуса $R=20$\,см, а захват требует попадания в шар радиуса $r=2$\,см в $\R^3$. Тогда одно событие имеет вероятность $(r/R)^3=10^{-3}$, а цепочка «захват--подъём--перенос--отпускание» из четырёх таких событий --- порядка $10^{-12}$. В задаче «кодового замка» нижняя граница экспоненциальна для \emph{любого} алгоритма без априорной информации~\cite{kakade-thesis}. Следовательно, декомпозиция на навыки не эвристика, а способ обойти неустранимый барьер.

\begin{proposition}[Плотная награда: сигнал и мост к бинарной]\label{prop:dense}
\begin{enumerate}[label=(\alph*)]
\item Если условный закон плотной отдачи $R_{\rm dense}$ при данном контексте не имеет атомов, то $\Pmix=1$: каждая группа информативна.
\item Пусть успех содержит событие $\{\min_hd(s_h)\le r\}$, где $d$ --- расстояние до цели, и $g\ge0$ неубывает, $g(r)>0$. Если плотная стадия добилась $\E\bigl[\tfrac1H\sum_hg(d(s_h))\bigr]\le\bar g$, то $J\ge1-\bar g/g(r)$. Для конечной награды $\E g(d(s_H))\le\bar g$ та же оценка верна для события $\{d(s_H)\le r\}$.
\item Максимум плотной цели может иметь нулевую вероятность успеха.
\end{enumerate}
\end{proposition}
\begin{proof}
(а) Для независимых безатомных величин $\PP(R_1=R_2)=0$.
(б) Если $d(s_h)>r$ при всех $h$, то $\tfrac1H\sum_hg(d(s_h))\ge g(r)$. Остаётся применить неравенство Маркова.
(в) Пусть успех есть $\{d(s_H)\le r\}$, а награда равна $-d(s_H)$. Политика A всегда заканчивает на расстоянии $r+\epsilon$: её плотная цель $-(r+\epsilon)$, а успех $0$. Политика B с вероятностью $\tfrac12$ попадает точно, иначе заканчивает на расстоянии $3r$: её плотная цель $-1{,}5r<-(r+\epsilon)$ при $\epsilon<r/2$, а успех $\tfrac12$.
\end{proof}

Отсюда порядок «плотная стадия $\to$ бинарная стадия». Плотная награда даёт сигнал с нулевой стартовой вероятности успеха (п.~(а)), но оптимизирует смещённую цель (п.~(в)). Если она довела ожидаемое расстояние до уровня ниже радиуса успеха, то по п.~(б) бинарная стадия начинается с $p_0>0$ и далее оптимизирует саму $J$ (теорема~\ref{thm:arcsin}). Требование «окна плотной награды шире радиуса успеха» есть требование, чтобы $g$ различала состояния вне успеха.

\begin{theorem}[Лестница стартовых распределений]\label{thm:ladder}
Пусть $\mu_0,\dots,\mu_N$ --- ступени (от лёгких стартов к целевому распределению). Переход на ступень $j+1$ выполняется, когда оценка успеха на ступени $j$ достигает порога $P$; возврат --- когда она падает ниже $D<P$. Допустим, что в момент перехода
\begin{equation}\label{eq:drop}
J_{\mu_{j+1}}(\pi)\ge J_{\mu_j}(\pi)-\Delta .
\end{equation}
Достаточные условия: $\Delta\ge\TV(\mu_j,\mu_{j+1})$, или $\Delta\ge L_VW_1(\mu_j,\mu_{j+1})$, если $V^\pi_0$ $L_V$-липшицева по стартовому состоянию. Тогда:
\begin{enumerate}[label=(\roman*)]
\item если $P-\Delta\ge D$ (\emph{условие храповика}), то сразу после каждого перехода успех не ниже $D$ и немедленного возврата нет;
\item пока успех на новой ступени остаётся в $[P-\Delta,P]$, доля смешанных групп не меньше $\min\{\Pmix(P-\Delta,G),\Pmix(P,G)\}$, то есть сигнал не исчезает;
\item если стадия RL, начатая с успеха не ниже $s$, достигает $P$ за $T(s)$ обновлений, то весь путь занимает не более $N\,T(P-\Delta)$ обновлений;
\item условие (i) точно: при $P-\Delta<D$ существуют допустимые ступени, после перехода на которые сразу срабатывает возврат.
\end{enumerate}
При оценке успеха по $n$ эпизодам с точностью $\zeta_n=\sqrt{\log(2/\delta)/(2n)}$ (Хёфдинг) условие принимает вид $P-\Delta-2\zeta_n\ge D$.
\end{theorem}
\begin{proof}
Достаточные условия: $J_{\mu'}-J_\mu=\int V^\pi_0\,d(\mu'-\mu)\ge-\TV(\mu,\mu')$, так как $V\in[0,1]$, и $\ge-L_VW_1$ по двойственности Канторовича--Рубинштейна.
(i) $J_{\mu_{j+1}}\ge P-\Delta\ge D$.
(ii) Функция $\Pmix(\cdot,G)$ унимодальна с максимумом в $\tfrac12$, поэтому минимум на отрезке достигается на конце.
(iii) Суммирование по ступеням.
(iv) Пример с равенством в~\eqref{eq:drop}.
Поправка на оценивание: на событии $|\hat J-J|\le\zeta_n$ для обеих оценок.
\end{proof}

Теорема формализует правило, найденное в экспериментах CDPR (раздел~\ref{sec:experiments}): старые пороги $0{,}30/0{,}20$ при измеренном $\Delta_{p90}=0{,}129$ нарушают условие ($0{,}171<0{,}20$), новые $0{,}45/0{,}30$ удовлетворяют ему ($0{,}321\ge0{,}30$). При $P-\Delta=0{,}321$ и $G=8$ имеем $\Pmix\approx0{,}955$.

\begin{theorem}[Композиция навыков под одной инструкцией]\label{thm:composition}
Пусть политика реактивна (замечание~\ref{rem:reactive}), а составная задача $\cT_c=(\ell_c,\mu_c,E_c)$ состоит из этапов $j=1,\dots,m$ с моментами остановки $0=\tau_0<\tau_1<\dots<\tau_m$ и событиями «этап $j$ завершён». Для каждого этапа задан навык $\cT_j=(\ell_j,\mu_j,E_j)$ с тем же событием завершения этапа. Пусть $V^\ell_j(s)$ --- вероятность завершить этап $j$ из состояния $s$ при инструкции $\ell$ (при достаточном оставшемся горизонте). Обозначим:
\begin{itemize}
\item $p_j=\E_{\mu_j}V^{\ell_j}_j$ --- качество навыка;
\item $\rho_j$ --- закон состояния передачи $s_{\tau_j}$ при условии завершения этапов $1..j$ под инструкцией $\ell_c$, $\rho_0=\mu_c$;
\item $\Gamma_j=\sum_h\E_{d^{\pi,\cT_j}_h}\TV\bigl(\pi(\cdot\mid s,\ell_c),\pi(\cdot\mid s,\ell_j)\bigr)$ --- разрыв переноса инструкции.
\end{itemize}
Тогда
\[
J_{\cT_c}(\pi)=\prod_{j=1}^m\E_{\rho_{j-1}}V^{\ell_c}_j\ \ge\ \prod_{j=1}^m\Bigl(p_j-\Gamma_j-\TV(\rho_{j-1},\mu_j)\Bigr)^+ .
\]
\end{theorem}
\begin{proof}
По строгому марковскому свойству и реактивности условная вероятность завершить этап $j$ при условии завершения этапов $1..j-1$ и при данном $s_{\tau_{j-1}}$ равна $V^{\ell_c}_j(s_{\tau_{j-1}})$. Поэтому
$\PP(E^{(j)})=\PP(E^{(j-1)})\,\E_{\rho_{j-1}}V^{\ell_c}_j$, откуда следует равенство. Далее, $V\in[0,1]$, поэтому $\E_{\rho}V^{\ell_c}_j\ge\E_{\mu_j}V^{\ell_c}_j-\TV(\rho,\mu_j)$. Теорема~\ref{thm:retention}, применённая к задаче $(\mu_j,E_j)$ и паре политик $\pi(\cdot\mid\cdot,\ell_c)$, $\pi(\cdot\mid\cdot,\ell_j)$, даёт $\E_{\mu_j}V^{\ell_c}_j\ge p_j-\Gamma_j$. Множители неотрицательны, поэтому можно взять положительные части.
\end{proof}

Каждый член этой оценки соответствует механизму пайплайна и измерим.
\begin{itemize}
\item $p_j$ --- качество навыков; его обеспечивают стадии RL навыков.
\item $\Gamma_j$ --- насколько политика под составной инструкцией ведёт себя так же, как под инструкцией навыка. Для гауссовской головы $\Gamma_j\le\sum_h\E\sqrt{\tfrac12\sum_i(\mu_i(x,\ell_c)-\mu_i(x,\ell_j))^2/\sigma_i^2}$; это вычисляется по записанным состояниям без новых эпизодов. Переразметка собственных успехов навыков под итоговую инструкцию (composition SFT) есть прямая минимизация $\Gamma_j$. Если навык размещения обучается под самой итоговой инструкцией, $\Gamma=0$ для этого этапа.
\item $\TV(\rho_{j-1},\mu_j)$ --- рассогласование передачи: насколько старты навыка (например, «объект в захвате») покрывают реальные состояния после предыдущего этапа. Его можно оценить двухвыборочным классификатором.
\end{itemize}
Сравнение с предложением~\ref{prop:barrier} показывает выигрыш декомпозиции: без навыков $J_c(\pi_0)\le q^m$, с навыками $J_c(\pi_0)$ отделена от нуля величиной, не зависящей экспоненциально от $m$. Например, при $\mu_1=\mu_c$ (так что $\TV(\rho_0,\mu_1)=0$), $p_1=0{,}5$, $\Gamma_1=0{,}1$, $p_2=0{,}7$, $\TV(\rho_1,\mu_2)=0{,}3$, $\Gamma_2=0$ получаем $J_c\ge0{,}16$ и $\Pmix(0{,}16,8)\approx0{,}75$.

\begin{proposition}[Самоимитация как шаг улучшения]\label{prop:sil}
Пусть $\beta$ --- политика, порождающая банк успехов задачи $\cT$, $q^+_\beta$ --- закон её траекторий при условии успеха, $J_\beta>0$. Для любой $\theta$
\[
\log\frac{J_\cT(\pi_\theta)}{J_\cT(\pi_\beta)}\ \ge\ G_\cT(\theta):=\E_{q^+_\beta}\sum_h\log\frac{\pi_\theta(u_h\mid x_h)}{\pi_\beta(u_h\mid x_h)} .
\]
В частности, если подгонка по правдоподобию на собственных успехах не уменьшила популяционное правдоподобие ($G_\cT\ge0$), то $J_\cT(\pi_\theta)\ge J_\cT(\pi_\beta)$.
\end{proposition}
\begin{proof}
Аналогично следствию~\ref{cor:anchor-success}: $J_\theta\ge J_\beta\E_{q^+}[d\PP^\theta/d\PP^\beta]$, затем неравенство Йенсена; множители динамики сокращаются~\cite[\S6]{notes}.
\end{proof}

Это шаг типа EM для RL~\cite{dayan-hinton,rwr}. Самоимитация на собственных успехах является оператором улучшения политики, а не внешним надзором. Для гауссовской головы с фиксированной ковариацией логарифм правдоподобия есть взвешенная MSE \emph{на латентном} $u$. Подгонка к усечённым или сглаженным действиям этой гарантии формально не наследует.

\begin{corollary}[Сбалансированная самоимитация задаёт нижнюю границу каждого навыка]\label{cor:sil-floor}
Пусть подгонка ведётся на смеси банков с весами $w_\cT>0$ и $\sum_\cT w_\cT G_\cT(\theta)\ge0$. Тогда $\sum_\cT w_\cT\log(J_{\theta,\cT}/J_{\beta,\cT})\ge0$, и для каждой задачи
\[
J_{\theta,\cT}\ \ge\ J_{\beta,\cT}\prod_{\cT'\ne\cT}J_{\beta,\cT'}^{\,w_{\cT'}/w_\cT} .
\]
Задача с $w_\cT=0$ нижней границы не имеет.
\end{corollary}
\begin{proof}
Просуммировать неравенства предложения~\ref{prop:sil} с весами. Затем воспользоваться тем, что $\log(J_{\theta,\cT'}/J_{\beta,\cT'})\le\log(1/J_{\beta,\cT'})$ при $J\le1$.
\end{proof}

Сбалансированная самоимитация поднимается по взвешенному среднему геометрическому $\prod_\cT J_\cT^{w_\cT}$, которое обращается в нуль при исчезновении любого навыка. Это логарифмический барьер против забывания, отсутствующий у RL на одном навыке (следствие~\ref{cor:mixture}).

%======================================================================
\section{Сводная теорема: ответ на вопрос}\label{sec:main}

\begin{theorem}[О бутстрэппинге исполнительной части VLA-политики]\label{thm:main}
Рассмотрим пайплайн: (I) стадии навыков $\cT_1,\dots,\cT_m$ (плотная, затем бинарная награда, с самоимитацией); (II) переход к составной инструкции; (III) совместное RL на $\cT_c$ (и, возможно, на навыках из смеси) из политики $\pi^{\rm III}_0$. Пусть выполнены допущение~\ref{ass:features} и следующие условия.
\begin{enumerate}[label=(A\arabic*)]
\item \emph{Представимость.} Существует $\pit\in\Pi_\varphi$ с $\bar D:=D(\pit\|\pi^{\rm III}_0)<\infty$.
\item \emph{Сигнал.} После стадий (I)--(II) выполнено $p_0:=\prod_j\bigl(p_j-\Gamma_j-\TV(\rho_{j-1},\mu_j)\bigr)^+>0$ (теорема~\ref{thm:composition}).
\item \emph{Неточность шага.} $\sum_h\sqrt{C_h}\,\|e^+_{k,h}\|_{L^2(\nu_h)}\le\varepsilon_{\rm upd}$ при всех $k$ (теорема~\ref{thm:md-inexact}).
\end{enumerate}
Тогда:
\begin{enumerate}[label=(\alph*)]
\item вероятность того, что группа на первом обновлении (III) содержит хотя бы один успех, не меньше $p_0$, а смешанной группы --- не меньше $p_0-\E_cp_c^G$; пока шаг нулевой, политика не меняется, поэтому ожидаемое число групп до первого ненулевого шага не больше $1/(p_0-\E_cp_c^G)$, если знаменатель положителен;
\item при $\eta=\sqrt{8\bar D/(HK)}$
\[
\max_{k<K}J_{\cT_c}(\pi_k)\ \ge\ J_{\cT_c}(\pit)-\sqrt{\frac{H\bar D}{2K}}-\varepsilon_{\rm upd};
\]
\item для навыка с весом $w_\cT>0$ в смеси шаг с дрейфом $\bar\Delta_k$ меняет $J_\cT$ не более чем на $\bar\Delta_k/w_\cT$; с якорём бюджета $D_{\rm ref}$ отклонение от эталона не превосходит $\sqrt{HD_{\rm ref}/2}$ на всех итерациях;
\item $J_{\cT_c}(\pit)\le J^\ast_\varphi$, и разрыв до $\sup J$ определяется представлением (предложение~\ref{prop:ceiling}).
\end{enumerate}
\end{theorem}
\begin{proof}
(а) При случайном контексте вероятность хотя бы одного успеха в группе равна $\E_c[1-(1-p_c)^G]\ge\E_cp_c=J_{\cT_c}(\pi^{\rm III}_0)\ge p_0$ по теореме~\ref{thm:composition}; вычитая вероятность $\E_cp_c^G$ группы из одних успехов, получаем оценку для смешанной группы. До первого ненулевого шага политика неподвижна (лемма~\ref{lem:stationary}), так что группы независимы и одинаково распределены.
(б) Теорема~\ref{thm:md-inexact} с выбранным $\eta$; среднее не превосходит максимума.
(в) Следствия~\ref{cor:mixture} и~\ref{cor:anchor}.
(г) Определение $J^\ast_\varphi$ и $\pit\in\Pi_\varphi$.
\end{proof}

\begin{proposition}[Необходимость сигнала]\label{prop:necessity}
Если $p_c(\pi^{\rm III}_0)=0$ для $\cD$-п.в.\ $c$ и регуляризация отсутствует, то $\pi_k=\pi^{\rm III}_0$ при всех $k$. Тогда оценка (б) верна лишь с $\varepsilon_{\rm upd}\ge J(\pit)-J(\pi^{\rm III}_0)-\sqrt{H\bar D/(2K)}$, то есть условие~(A3) нарушается. Сигнал~(A2) необходим для~(A3), а простые навыки --- конструктивный способ его обеспечить.
\end{proposition}
\begin{proof}
Лемма~\ref{lem:stationary}: все группы вырождены почти наверное, шаг нулевой, и неравенство (б) превращается в утверждение о $\varepsilon_{\rm upd}$.
\end{proof}

\paragraph{Ответ на вопрос.} Исполнительную часть VLA-политики можно обучить только консервативным RL по четырём причинам.
(У1)~При бинарной награде GRPO поднимается по монотонному преобразованию вероятности успеха, то есть оптимизирует саму метрику задачи (теорема~\ref{thm:arcsin}).
(У2)~Шаг в доверительной области почти монотонен и меняет успех каждого навыка, на состояниях которого он ограничен, не более чем на сумму средних вариаций (теоремы~\ref{thm:improve}, \ref{thm:retention}).
(У3)~Последовательность таких шагов есть неточный зеркальный спуск, сходящийся к любому успешному поведению, представимому на замороженных признаках. Цена определяется KL-расстоянием от начальной политики до этого поведения, покрытием и невязкой шага (теоремы~\ref{thm:md}, \ref{thm:md-inexact}).
(У4)~Сигнал существует, только если начальная вероятность успеха отделена от нуля. Для составной задачи без декомпозиции она экспоненциально мала (предложение~\ref{prop:barrier}), а описание простых навыков тройками (инструкция, старты, предикат успеха) с плотной наградой, лестницей стартов и композицией делает её отделённой от нуля за полиномиальную цену (предложение~\ref{prop:dense}, теоремы~\ref{thm:ladder}, \ref{thm:composition}).

\paragraph{Что дали бы демонстрации и чем пайплайн их заменяет.} В сводной оценке демонстрации могут войти лишь через две величины:
через $\bar D$ (клонирование поведения приближает $\pi_0$ к $\pit$)
и через $C_h$ (рестарты из состояний демонстраций обеспечивают покрытие~\cite{salimans-chen}).
Пайплайн управляет $\bar D$ с помощью RL навыков и самоимитации на собственных успехах (предложение~\ref{prop:sil}), а $C_h$ --- с помощью стартовых распределений навыков, доступных в симуляторе (замечание~\ref{rem:reactive}). Поэтому оценка не требует демонстраций. В этом точный математический смысл утверждения «датасет для новой структуры не обязателен». При развёртывании привилегированные старты не используются: итоговая оценка проводится из $\mu_c$.

%======================================================================
\section{Связь с экспериментами и проверяемые предсказания}\label{sec:experiments}

\begin{longtable}{>{\raggedright}p{0.24\linewidth}>{\raggedright}p{0.30\linewidth}>{\raggedright\arraybackslash}p{0.38\linewidth}}
\toprule
Величина теории & Механизм и диагностика & Наблюдение \\
\midrule\endhead
$\Pmix$, вырожденные группы (лемма~\ref{lem:stationary}, зам.~\ref{rem:constants}) & динамическое сэмплирование, фильтр std $0{,}05$; доля смешанных групп по стадиям & фильтр при бинарной награде удаляет ровно вырожденные группы \\
$\ab_h\le\sqrt{\overline{\KL}/2}$ (лемма~\ref{lem:pinsker}) & усечение $0{,}2/0{,}28$, отбраковка расходящихся шагов; \texttt{approx\_kl\_mean} & CDPR, LR $10^{-5}$: $\overline{\KL}=0{,}009$ на обновление, $\ab\le0{,}067$ \\
Накопление дрейфа (сл.~\ref{cor:drift}) & сумма $\sum_k\sqrt{\overline{\KL}_k/2}$ по запуску & CDPR: за 113 обновлений при малом KL на шаге потеряно удержание подъёма; «обмен» plate/bowl и захват/размещение \\
Якорь (сл.~\ref{cor:anchor}, \ref{cor:anchor-success}) & \texttt{reference\_anchor}, член $(\mu-\mu_{\rm ref})^2/(2\sigma^2_{\rm ref})$ на банке успехов & вдвое меньший ущерб plate/lift, но эрозия lift$\mid$grasp сохраняется: граница зависит от бюджета $D_{\rm ref}$ и от того, покрывает ли банк нужные состояния \\
$w_\cT=0$ (сл.~\ref{cor:mixture}, \ref{cor:sil-floor}) & стадии RL по отдельным инструкциям; поинструкционное удержание & WidowX: подход и захват потеряны уже на стадии размещения (s3); CDPR: захват стёрт до 0, когда отсутствовал в смеси SFT \\
$D(\pit\|\pi_0)$ (сл.~\ref{cor:gauss}) & остаток над приором; насыщение $\tanh$ & CDPR: приор --- постоянное смещение, остаток должен его компенсировать и насыщен \\
Потолок $J^\ast_\varphi$ (пр.~\ref{prop:ceiling}) & замороженные признаки; ошибка локализации против допуска & CDPR: локализация 3--5\,см при допуске захвата около 2\,см ограничивает захват; размещение с допуском 5,7/9,1\,см устойчиво \\
Храповик $P-\Delta\ge D$ (т.~\ref{thm:ladder}) & пороги лестницы; $\Delta_{p90}$ & $0{,}30/0{,}20$: 3 из 4 семейств «застряли»; $0{,}45/0{,}30$: ноль возвратов, захват выведен из нуля \\
$\Gamma_j$, $\TV(\rho,\mu)$ (т.~\ref{thm:composition}) & переразметка успехов под итоговую инструкцию; старты «в захвате» & CDPR: составная задача отказывает на передаче захват$\to$отпускание; WidowX: переход composition SFT оказался неэффективным \\
\bottomrule
\end{longtable}

\paragraph{Предсказания для эксперимента WidowX-200.} Теория даёт конкретные измерения, подтверждающие или опровергающие механизм.
\begin{enumerate}
\item \emph{Сигнал.} Совместное RL должно стартовать, только когда оценка $J_c(\pi^{\rm III}_0)$ заметно больше нуля. Доля смешанных групп в начале (III) должна согласовываться с $\E_c\Pmix(p_c,8)$.
\item \emph{Перенос инструкции.} До и после composition SFT вычислить $\hat\Gamma_1$ по записанным состояниям навыка захвата (KL между средними головы при $\ell_c$ и $\ell_{\rm pick}$). Теорема~\ref{thm:composition} предсказывает: если $\hat\Gamma_1$ не уменьшилось, переход не поднимет $J_c$. Это согласуется с наблюдённой неэффективностью перехода.
\item \emph{Передача.} Оценить $\TV(\rho_1,\mu_2)$ между реальными состояниями после захвата и стартами «в захвате». Большое значение предсказывает отказы на передаче независимо от качества навыка размещения.
\item \emph{Сохранение.} Держать каждую инструкцию в смеси с $w_\cT\ge w_{\min}$ или под якорём. Сообщать накопленный бюджет $\sum_k\sqrt{\overline{\KL}_k/2}$ как верхнюю границу дрейфа и проверять, что каждая семья выше своей нижней границы (следствие~\ref{cor:sil-floor}).
\item \emph{Храповик.} Для каждой лестницы проверять $P-\Delta_{p90}-2\zeta_n\ge D$ до запуска.
\end{enumerate}

%======================================================================
\section{Границы применимости}\label{sec:limits}

\begin{enumerate}
\item \textbf{Усечение не является доверительной областью.} Теоремы раздела~\ref{sec:stability} относятся к \emph{реализованным} шагам. Выполнение их условий надо измерять, а не выводить из значения $\varepsilon_\pm$.
\item \textbf{Константы худшего случая.} При $H\sim40$--$128$ оценки улучшения численно бессодержательны (замечание~\ref{rem:clip}). Их роль --- масштабирование и структура. Теорема~\ref{thm:retention} численно почти точна, но требует TV по всем шагам задачи.
\item \textbf{Неточная сходимость --- разложение, а не гарантия.} Для нейросетевой политики невязка $e_k$ теорией не ограничивается. Теорема~\ref{thm:md-inexact} говорит, \emph{куда} входят оценивание, представление и покрытие, но не доказывает, что они малы.
\item \textbf{Обучение против развёртывания.} Цель $J$ относится к стохастической обучающей политике. Оценка CDPR использует детерминированное среднее остатка при стохастическом приоре; $J_{\rm explore}$ и $J_{\rm deploy}$ надо измерять отдельно~\cite{notes}.
\item \textbf{Реактивность.} Теоремы о композиции и о рестартах используют реактивность политики (замечание~\ref{rem:reactive}); память кусков действий требует её включения в состояние.
\item \textbf{Происхождение данных.} Аудит происхождения отчётного контроллера CDPR (\path{ral_submission/evidence/DATA_LINEAGE_AUDIT_20261029.md}) показывает, что человеческих демонстраций нет. Однако в предках есть стадии SFT на данных скриптового оракула и на эпизодах выученной политики с помощью скриптовых контроллеров. Теория объясняет, почему такие стадии помогали: они дают $p_0>0$ и малое $\bar D$. Но для результата 69{,}14\,\% утверждение «без каких-либо подготовленных данных» не подтверждается. Производственная конфигурация WidowX использует только данные, порождённые самой политикой (оракул --- лишь в дымовом тесте), и потому является чистой проверкой теоремы~\ref{thm:main}.
\item \textbf{Атрибуция.} Лемма~\ref{lem:pdl} восходит к~\cite{kakade-langford}. Теорема~\ref{thm:improve} --- конечногоризонтная версия рассуждения~\cite{trpo,cpo}. Теорема~\ref{thm:md} --- адаптация анализа зеркального спуска в пространстве политик~\cite{shani,lan,xiao,agarwal}. Теорема~\ref{thm:arcsin}(а) --- известное свойство нормированного GRPO~\cite{mroueh}; см.\ также~\cite{dr-grpo} о смещении нормировки. Предложение~\ref{prop:sil} --- неравенство EM-типа~\cite{dayan-hinton,rwr}. Предложение~\ref{prop:barrier} --- вариант нижних границ «кодового замка»~\cite{kakade-thesis}. Собственными для данной работы являются формулировки для вероятности успеха с явными константами, следствия~\ref{cor:mixture}--\ref{cor:anchor-success} и~\ref{cor:sil-floor}, теоремы~\ref{thm:ladder}, \ref{thm:composition}, \ref{thm:main} и их связь с измеримыми величинами пайплайна. Окончательную новизну надо подтвердить обзором литературы.
\end{enumerate}

%======================================================================
\section{Проект положений, выносимых на защиту}\label{sec:defense}

\begin{enumerate}
\item Для эпизодических задач манипуляции с бинарным предикатом успеха групповой метод GRPO в точке поведения является несмещённым подъёмом по строго возрастающему преобразованию вероятности успеха каждого контекста. Его стационарные точки Парето-оптимальны по вектору вероятностей успеха контекстов.
\item Для вероятности успеха на конечном горизонте $H$ получены оценки улучшения в доверительной области с константой $H(H-1)/2$ и оценка сохранения навыка через сумму средних вариаций политики на состояниях этого навыка. Из них следуют количественные условия сохранения навыков при обучении смеси инструкций, линейное накопление дрейфа при покомпонентно малых шагах и равномерная по времени граница для якоря к эталонной политике.
\item Консервативная итерация политики с групповыми монте-карловскими оценками трактуется как неточный зеркальный спуск. Для неё доказана оценка разрыва до любой сравнимой политики над замороженными признаками, разложенная на оптимизационный член $O(\sqrt{HD/K})$, член покрытия ($\chi^2$-концентрируемость) и невязку шага.
\item Доказано, что при разреженной награде начальная вероятность успеха необходима для сигнала, и для составной задачи она экспоненциально мала по числу этапов. Спецификация навыков (инструкция, стартовое распределение, предикат успеха) вместе с плотной наградой, лестницей стартов с условием храповика $P-\Delta\ge D$ и композицией обеспечивает её нижнюю границу через качество навыков, разрыв переноса инструкции и рассогласование передачи.
\item В полученной оценке демонстрации участвуют только через KL-расстояние начальной политики до успешной и через коэффициент покрытия. Обе величины управляемы без демонстраций средствами симулятора и самоимитации, что является математическим обоснованием адаптации VLA-моделей к новым кинематическим структурам без ручного сбора данных.
\end{enumerate}

%======================================================================
\appendix
\section{Численная проверка утверждений}\label{app:numerics}

Скрипт \texttt{verify\_theorems.py} проверяет неравенства на случайных конечных марковских процессах с конечным горизонтом и терминальной наградой успеха. Величины $J$, $Q$, $A$ и $d_h$ вычисляются точно динамическим программированием. Результаты записываются в \texttt{verification\_results.json}. Проверка ловит ошибки в алгебре и константах, но не сертифицирует нейросетевую политику.

\begin{center}
\begin{tabularx}{\linewidth}{>{\raggedright}p{0.36\linewidth}r>{\raggedright\arraybackslash}X}
\toprule
Утверждение & Случаев & Результат \\
\midrule
Лемма~\ref{lem:pdl} (тождество) & 4000 & ошибка $\le6\cdot10^{-16}$ \\
Лемма~\ref{lem:shift} & 4000 & выполнено \\
Теорема~\ref{thm:improve}, все три формы & 4000 & выполнено; наибольшее отношение «нарушение/штраф» $0{,}56$ \\
Теорема~\ref{thm:retention} & 4000 & выполнено; наибольшее отношение $|\Delta J|/\Delta_\cT=0{,}987$ (почти точна) \\
Следствие~\ref{cor:anchor} & 4000 & выполнено \\
Теорема~\ref{thm:md} (точный шаг Гиббса) & 600 & выполнено; наибольшее отношение разрыва к оценке $0{,}71$ \\
Теорема~\ref{thm:arcsin}(б), $p\in\{0{,}03..0{,}8\}$, $G\in\{2,8,32\}$ & $4\cdot10^5$ групп & $|z|\le2{,}26$ (Монте-Карло против формулы) \\
Теорема~\ref{thm:composition} (шаг передачи) & 20000 & выполнено \\
\bottomrule
\end{tabularx}
\end{center}

Запуск из каталога главы: \texttt{python3 verify\_theorems.py}.

%======================================================================
\begin{thebibliography}{99}\small
\bibitem{notes} Mathematical foundation for the RA-L paper. Рабочая заметка, \path{ral_submission/theory/MATHEMATICAL_FOUNDATION.md}, 2026.
\bibitem{kakade-langford} Kakade S., Langford J. Approximately optimal approximate reinforcement learning // ICML. 2002.
\bibitem{kakade-thesis} Kakade S. On the sample complexity of reinforcement learning. PhD thesis, University College London, 2003.
\bibitem{trpo} Schulman J., Levine S., Abbeel P., Jordan M., Moritz P. Trust region policy optimization // ICML. 2015.
\bibitem{ppo} Schulman J., Wolski F., Dhariwal P., Radford A., Klimov O. Proximal policy optimization algorithms // arXiv:1707.06347. 2017.
\bibitem{cpo} Achiam J., Held D., Tamar A., Abbeel P. Constrained policy optimization // ICML. 2017.
\bibitem{grpo} Shao Z. et al. DeepSeekMath: Pushing the limits of mathematical reasoning in open language models // arXiv:2402.03300. 2024.
\bibitem{dapo} Yu Q. et al. DAPO: An open-source LLM reinforcement learning system at scale // arXiv:2503.14476. 2025.
\bibitem{mroueh} Mroueh Y. Reinforcement learning with verifiable rewards: GRPO's effective loss, dynamics, and success amplification // arXiv:2503.06639. 2025.
\bibitem{dr-grpo} Liu Z. et al. Understanding R1-Zero-like training: A critical perspective // arXiv:2503.20783. 2025.
\bibitem{engstrom} Engstrom L. et al. Implementation matters in deep RL: A case study on PPO and TRPO // ICLR. 2020.
\bibitem{truly-ppo} Wang Y., He H., Tan X. Truly proximal policy optimization // UAI. 2020.
\bibitem{agarwal} Agarwal A., Kakade S., Lee J., Mahajan G. On the theory of policy gradient methods: Optimality, approximation, and distribution shift // JMLR. 2021. Vol.~22, no.~98.
\bibitem{shani} Shani L., Efroni Y., Mannor S. Adaptive trust region policy optimization: Global convergence and faster rates for regularized MDPs // AAAI. 2020.
\bibitem{lan} Lan G. Policy mirror descent for reinforcement learning: Linear convergence, new sampling complexity, and generalized problem classes // Mathematical Programming. 2023. Vol.~198.
\bibitem{xiao} Xiao L. On the convergence rates of policy gradient methods // JMLR. 2022. Vol.~23.
\bibitem{mdpo} Tomar M., Shani L., Efroni Y., Ghavamzadeh M. Mirror descent policy optimization // ICLR. 2022.
\bibitem{neural-ppo} Liu B., Cai Q., Yang Z., Wang Z. Neural trust region/proximal policy optimization attains globally optimal policy // NeurIPS. 2019.
\bibitem{dayan-hinton} Dayan P., Hinton G. Using expectation-maximization for reinforcement learning // Neural Computation. 1997. Vol.~9, no.~2.
\bibitem{rwr} Peters J., Schaal S. Reinforcement learning by reward-weighted regression for operational space control // ICML. 2007.
\bibitem{salimans-chen} Salimans T., Chen R. Learning Montezuma's Revenge from a single demonstration // arXiv:1812.03381. 2018.
\bibitem{jl} Johnson W., Lindenstrauss J. Extensions of Lipschitz mappings into a Hilbert space // Contemporary Mathematics. 1984. Vol.~26.
\bibitem{dasgupta} Dasgupta S., Gupta A. An elementary proof of a theorem of Johnson and Lindenstrauss // Random Structures \& Algorithms. 2003. Vol.~22, no.~1.
\bibitem{smolvla} Shukor M. et al. SmolVLA: A vision-language-action model for affordable and efficient robotics // arXiv:2506.01844. 2025.
\end{thebibliography}

\end{document}
